\documentclass[12pt]{article}
\usepackage{amsmath,amssymb,amsfonts}
\usepackage[margin=1in]{geometry}

\begin{document}

\textbf{Chapter 6, Problem 14.} \textit{The mean-value theorem.} Prove that for charge-free two-dimensional space the value of the electrostatic potential at any point is equal to the average of the potential over the surface of any circle centered on that point. Do this by considering the electrostatic potential as the real part of an analytic function.

\bigskip
\textbf{Solution.}

In a charge-free region, the electrostatic potential $\phi(x,y)$ satisfies Laplace's equation $\nabla^2\phi = 0$, so $\phi$ is harmonic. Since the region is simply connected, $\phi$ is the real part of some analytic function $f(z) = \phi + i\psi$.

By Cauchy's integral formula, for any point $z_0$ inside a circle $C$ of radius $R$ centered at $z_0$:
\[
f(z_0) = \frac{1}{2\pi i}\oint_C \frac{f(z)}{z-z_0}\,dz.
\]

Parametrize the circle: $z = z_0 + Re^{i\theta}$, $dz = iRe^{i\theta}\,d\theta$, $z - z_0 = Re^{i\theta}$. Then:
\[
f(z_0) = \frac{1}{2\pi i}\int_0^{2\pi}\frac{f(z_0+Re^{i\theta})}{Re^{i\theta}}\,iRe^{i\theta}\,d\theta = \frac{1}{2\pi}\int_0^{2\pi}f(z_0+Re^{i\theta})\,d\theta.
\]

Taking the real part of both sides:
\[
\phi(z_0) = \text{Re}\,f(z_0) = \frac{1}{2\pi}\int_0^{2\pi}\text{Re}\,f(z_0+Re^{i\theta})\,d\theta = \frac{1}{2\pi}\int_0^{2\pi}\phi(z_0+Re^{i\theta})\,d\theta.
\]

The right-hand side is precisely the average of $\phi$ over the circle of radius $R$ centered at $z_0$.

Therefore, the value of the potential at any point equals the average of the potential on any circle centered at that point (provided the circle lies within the charge-free region).

\end{document}
